\section{结论}
模型中的四大家鱼恢复到禁渔前约1.8倍，底层资源末值却降至0.204，食压指数持续上升。恢复主要积聚在消费者层，资源供给并未同步增加；这些数值仅表示归一化单区系统中的相对压力，不对应全流域实测生物量结构。恢复收益仍能传到长江江豚和长江鲟，只是这条路径并不稳定：通道阻隔会同时削弱两者，人工放流的补偿更多落在长江鲟。

承载力提高后，三级食物链不会一直维持规则振荡，长期波动随之复杂。有限时间最大李雅普诺夫指数支持这一变化方向，但还不足以确定唯一临界点；$K$也只表示归一化结构压力，而非长江某一河段的实测承载量。污染和入侵情景中，食源过程量可能与珍稀物种终态分离。禁渔只是恢复链条的起点；现有计算不能给出治理措施的实际效果，也不能由此断定恢复收益会长期留在系统内部。是否留得住，还取决于水质、连通性、放流和入侵控制。

